\section{Exclusion of all rank-one Lie cores}
\label{sec:a1-core}

The small-field part of this section uses complete published ATLAS
character tables.  The finite spectral statements have been checked
both by GAP and by a separate exact Ramanujan-trace implementation.
The outer-automorphism coverage requires a further argument at
$q=9,25$: a universal-cover extension missing from the table list
is not silently assumed to exist.

\subsection*{D-BI-21a. The complete small-field spectral input}

Let
\[
 \mathcal Q_0=\{4,5,7,8,9,11,13,16,17,19,23,25,27\}.
\]
The duplicate simple groups at $q=4,5$ are represented once.
For each simple group $\operatorname{PSL}_2(q)$, the following
table specifies its universal cover and all cyclic outer subgroups.
The notation after a semicolon indicates that a proper central
quotient is used for that outer subgroup, subject to D-BI-21b below.
All identifiers are exact CTblLib names.
\[
\begin{array}{c|l|l}
q&U&\text{cyclic outer extension tables}\\\hline
4,5&2.A5&2.A5.2\\
7&2.L3(2)&2.L3(2).2\\
8&L2(8)&L2(8).3\\
9&6.A6&6.A6.2\_1,\ 6.A6.2\_2;\quad3.A6.2\_3\\
11&2.L2(11)&2.L2(11).2\\
13&2.L2(13)&2.L2(13).2\\
16&L2(16)&L2(16).2,\ L2(16).4\\
17&2.L2(17)&2.L2(17).2\\
19&2.L2(19)&2.L2(19).2\\
23&2.L2(23)&2.L2(23).2\\
25&2.L2(25)&2.L2(25).2\_1,\ 2.L2(25).2\_2;
                    \quad L2(25).2\_3\\
27&2.L2(27)&2.L2(27).2,\ 2.L2(27).3,\ 2.L2(27).6
\end{array}                                                   \tag{BI.78}
\]
The two additional quotient-core tables $3.A6$ and $L2(25)$
are retained too.  Thus the full finite input has fourteen core
tables and nineteen extension tables, thirty-three in total.
Each extension is supplied with its actual stored class fusion
from the core named in the certificate.

\begin{lemma}[D-BI-21a: the finite spectral bounds]
\label{D-BI-21a}
In every nonlinear irreducible representation of a core table in
(BI.78), and in every irreducible representation of an extension
table whose restriction to its specified core is nonlinear and
irreducible, every nonscalar element satisfies
\[
                              \nu(g)\ge a/5,                 \tag{BI.79}
\]
where $a$ is the representation degree.  The stronger bound
$\nu(g)\ge a/4$ fails only for the following four entries:
\[
\begin{array}{c|r|r|r|r|r}
\text{table}&\text{row}&\text{class}&a&\operatorname{ord}(g)&\nu(g)\\\hline
6.A6.2\_1&3&20&5&2&1\\
6.A6.2\_1&4&20&5&2&1\\
6.A6.2\_1&5&21&5&4&1\\
6.A6.2\_1&6&21&5&4&1
\end{array}                                                   \tag{BI.80}
\]
The restrictions of these rows are, respectively, rows
$2,2,3,3$ of $6.A6$.  Those two core characters have degree five
and are trivial on the center.
\end{lemma}

\begin{proof}
The identifications, characters and power maps are explicit
published ATLAS inputs accessed through \cite{SPOR:CTblLib}.
Every retained table's source string identifies that origin.
The universal-cover and automorphism metadata are interpreted as
in the library's \S3.7-3, with the small multiplier and isomorphism
exceptions also listed in \cite[Chapter~5, Tables~1--2,
numbered pp.~35--36]{BI:HHM}.  No group identification is inferred
from character orthogonality alone.

For every selected character row $\chi$ and every class of an
element $g$ of order $m$, compute every eigenvalue multiplicity by
\[
 m_k=\frac1m\sum_{j=0}^{m-1}\chi(g^j)\zeta_m^{-kj}.
                                                               \tag{BI.81}
\]
All arithmetic is exact in a cyclotomic field.  The criterion
$|\chi(g)|^2=\chi(1)^2$ discards exactly the scalar classes.
An extension row is selected precisely when its restriction along
the actual core fusion has character norm one and degree greater
than one.  Hence no representation with an irreducible core
restriction is dropped by a degree cutoff or a reality condition.

The executed file \texttt{a1/export\_and\_scan.g} checks all
$654$ selected rows and $18\,922$ nonscalar entries.  It verifies
the integrality, nonnegativity and degree sum of every spectrum,
and produces exactly (BI.80) and no failure of (BI.79).  The full
input also includes every unselected character row, all class
sizes, group orders, power maps, source strings and subgroup
fusions.  Its SHA-256 is
\texttt{5ed905449e1f6f64f3eb18b27a213fc5c72e0731dbd591670166e8763ad7fc88}.

The separate program \texttt{check\_a1\_spectra.py} imports no GAP
calculation or selected-row list.  It selects the rows using exact
restriction norms and reconstructs all $18\,922$ spectra using
Ramanujan sums.  Explicitly, the normalized field trace of
$\zeta_m^j$ is $c_m(j)/\varphi(m)$, where
\[
                 c_m(j)=\sum_{d\mid\gcd(m,j)}d\mu(m/d).
\]
Since (BI.81) is a rational integer, its normalized field trace
equals itself.  This computes it independently by integer sums,
after the prime power maps and cyclotomic Galois action have
reconstructed the character values on all powers.  That program
finds the same four exceptions and verifies the stated core
restrictions and central kernels.  It also checks the complete
table-name inventory and all nineteen core fusions: orders,
identity kernels, full normal-class unions and compatibility
with every retained prime power map.  These complete finite checks
prove the lemma from the published input.  They do not reprove
the character tables or the classification of the named groups.
\end{proof}

\subsection*{D-BI-21b. The missing mixed outer extensions}

\begin{lemma}[D-BI-21b: a central obstruction]
\label{D-BI-21b}
Let $q=p^2$ with $p$ odd, and let $\alpha$ be the product of a
nontrivial diagonal outer automorphism and the field automorphism
of $\operatorname{SL}_2(q)$.  No irreducible complex representation
on which $-I$ acts as $-1$ is invariant under $\alpha$.
At $q=9$, the same assertion holds on the universal cover $6.A6$
for every representation nontrivial on its central involution.
\end{lemma}

\begin{proof}
Choose a generator $\zeta$ of $\mathbf F_q^*$ and take
\[
 \alpha(x)=D x^{(p)}D^{-1},\quad
 D=\operatorname{diag}(\zeta,1),\quad
 u=\operatorname{diag}(\zeta^{(p+1)/2},\zeta^{-(p+1)/2}).
\]
Direct matrix multiplication gives
\[
                     \alpha^2=\operatorname{inn}(u),
                     \qquad \alpha(u)=-u.                  \tag{BI.82}
\]
For the second equality, each diagonal exponent changes by
$\pm(p-1)(p+1)/2=\pm(q-1)/2$, whose power of $\zeta$ is $-1$.
These are the explicit diagonal and field automorphisms described
in \cite[\S3, pp.~5--7]{BI:WhiteSL2}.

Suppose an invertible intertwiner $T$ implements $\alpha$ in an
irreducible representation $\rho$.  The first equality in
(BI.82) and Schur's lemma imply $T^2=c\rho(u)$ for a nonzero
scalar $c$.  Since $T$ commutes with its own square, it must
commute with $\rho(u)$.  The second equality instead makes it
send $\rho(u)$ to $-\rho(u)$ when $\rho(-I)=-1$.
This contradiction proves the assertion for $\operatorname{SL}_2(q)$.

For $q=9$ let $U=6.A6$, let $U\to2.A6=\operatorname{SL}_2(9)$
be the canonical central quotient, and lift $\alpha$ uniquely
to $U$.  The uniqueness of this lift follows from the universal
property of the perfect central cover; two lifts with the same
action on $A_6$ differ by a homomorphism from the perfect group
$U$ to its center, and therefore coincide.  Choose any lift
$u'\in U$ of $u$.  Uniqueness gives
$\alpha^2=\operatorname{inn}(u')$ on $U$.
The element $z=\alpha(u')u'^{-1}$ is central, and its image in
$2.A6$ is $-I$.  Since $Z(U)$ is cyclic of order six, $z^3$ is
its central involution.  If that involution acts as $-1$, then
$\rho(z)\ne1$.  But an intertwiner would again satisfy
$T^2=c\rho(u')$ and commute with $\rho(u')$, while conjugating
it to $\rho(z)\rho(u')$.  This is impossible.
\end{proof}

\subsection*{D-BI-21c. Coverage of every linear normalizer}

\begin{lemma}[D-BI-21c: small-field normalizers]
\label{D-BI-21c}
Let $N$ be a finite quasisimple group with simple quotient
$\operatorname{PSL}_2(q)$ for $q\in\mathcal Q_0$, acting faithfully
and irreducibly on a space of dimension $a$.  Every nonscalar
element of its linear normalizer, with arbitrary scalar phase,
satisfies $\nu\ge a/5$.  It satisfies $\nu\ge a/4$ unless
the representation factors through $A_6$ and has degree five.
\end{lemma}

\begin{proof}
Lift the representation to the universal cover $U$.  The
normalizer modulo scalars embeds in $\operatorname{Aut}(S)$,
where $S=N/Z(N)$: a central automorphism of a perfect group is
trivial, and the linear centralizer is scalar.  The complete
outer group is generated by diagonal and field automorphisms,
with structure $C_{\gcd(2,q-1)}\times C_f$ for $q=p^f$;
see \cite[\S2, pp.~2--3; \S3, p.~6]{BI:WhiteSL2}.
For $A_5$ we use the equivalent $q=5$ model.  Thus (BI.78)
includes every cyclic subgroup generated by any possible outer
element, not only the largest outer group.  In particular both
orders two and four occur at $q=16$, all three subgroups of order
two occur at $q=9,25$, and orders two, three and six occur at $q=27$.

We spell out why a fixed cyclic extension table controls any
actual normalizing matrix.  Let $A$ be such a table group,
$A/U$ cyclic of order $l$, and choose $t\in A$ whose coset
generates it.  Write $\beta=\operatorname{inn}(t)|_U$ and
$u_0=t^l\in U$.  A normalizing matrix $g$, after an inner
adjustment and replacing $t$ by a suitable generating power,
induces the same automorphism of $S$ as $t$.  Unique lifting to
the universal perfect cover then gives
\[
                  g\rho(u)g^{-1}=\rho(\beta(u)).
\]
In particular the represented kernel is preserved.  Schur's
lemma makes $g^l\rho(u_0)^{-1}$ a nonzero scalar.  Rescale $g$
by an $l$th root to arrange $g^l=\rho(u_0)$.  The conjugation
and power relations now define an irreducible representation
of the actual fixed group $A$, extending $\rho$.  The original
matrix is, up to scalar and the inner adjustment, an element
of that represented extension.  Its spectrum is covered by
D-BI-21a.  This reasoning does not assume that $A/U$ splits,
nor that every core character is invariant under every outer
automorphism.

There are two exceptions to using the universal-cover extension
literally.  At $q=25$, the ATLAS subscripts $2_1,2_2,2_3$ denote
the diagonal, field, and mixed diagonal-field classes, respectively.
At $q=9$ the order of the first two is reversed: $2_1$ is the
field extension $S_6$, $2_2$ is the diagonal extension
$\operatorname{PGL}_2(9)$, and $2_3$ is the remaining mixed
extension.  The construction labels and exact subgroup fusions
are retained in the source inventory.  The missing usual cover
extensions are precisely these mixed cases.

D-BI-21b shows that an invariant core representation in either
mixed case must kill its central involution.  Hence it factors
through $L2(25)$ at $q=25$, and through $3.A6$ at $q=9$.
The actual extensions $L2(25).2_3$ and $3.A6.2_3$ and their
core fusions are present in (BI.78), so the preceding cyclic
argument applies to them.  No inference from an unavailable
universal-cover extension is made.  The resulting representations
are exactly among the rows tested in D-BI-21a.

Finally scalar multiplication does not change $\nu$, and a
normalizer element with inner outer action is scalar times an
element of $N$.  This proves the universal fifth bound and the
stated sole possible quarter-bound exception.
\end{proof}

\subsection*{D-BI-21. The rank-one core is impossible}

\begin{theorem}[D-BI-21: full $A_1$-core exclusion]
\label{D-BI-21}
Let a finite-cover Weil sixfold source of exact imaginary
quadratic character field have nonsplit actual polarization and
period rank nine.  Let an inducing block have degree $e\ge5$
and finite image $J$ that is primitive and tensor indecomposable
in the factorwise convention of D-BI-19.  Then the simple factors
of its forced normal core cannot be $\operatorname{PSL}_2(q)$
for any prime power $q$ for which this group is simple.

The statement includes the exceptional identifications $A_5$ and
$A_6$, the exceptional universal cover $6.A6$, arbitrary larger
inducing character fields, compatible outer automorphisms and
permutations of tensor factors.  For the small fields
$q\in\mathcal Q_0$, total degree $e\ge6$ is excluded even
without the nonsplit-polarization hypothesis.
\end{theorem}

\begin{proof}
D-BI-19 supplies a normal irreducibly acting central product of
quasisimple factors, permuted transitively by $J$.  Its factor
representations are faithful and irreducible, and their degrees
are equal, say $a$, so $e=a^t$.

If $q\ge29$, the published projective minimum-degree bound
$a\ge(q-1)/\gcd(2,q-1)$ in
\cite[Table~1, author p.~14]{BI:Hoffman} gives $e\ge14$.
The simple quotient is then nonalternating, so D-BI-20 applies
and excludes this case.  The remaining simple $A_1$ groups have
exactly the parameters in $\mathcal Q_0$.

For these parameters, every factorwise nonscalar normalizer
element satisfies $\nu\ge e/5$, by D-BI-21c and the tensor
eigenspace bound.  Every element permuting factors nontrivially
satisfies the stronger $\nu\ge e/4$, by the weighted cyclic-word
proof of D-SPOR-03; that proof does not depend on the factor
being sporadic.  Thus every nonscalar element of $J$ satisfies
the fifth bound.  Genuine scalar inertias have full moving rank.

The actual intermediate Hurwitz factorization gives genus zero,
at least twelve genuine branch points and total moving rank
$2e+6$, as proved in D-BI-18.  Consequently
\[
          12\lceil e/5\rceil\le2e+6,
          \qquad e\le15.                                  \tag{BI.83}
\]
This is a necessary constraint on the actual period map.
If no factor has the exceptional degree-five $A_6$ representation,
the identical argument uses the quarter bound and gives
\[
                         12\lceil e/4\rceil\le2e+6.
                                                               \tag{BI.84}
\]
For $e=5,6$ the left side is $24$; for $e\ge7$ it is at
least $3e>2e+6$.  Thus (BI.84) is impossible for every $e\ge5$.

If there is an exceptional factor, transitivity makes the common
factor degree $a=5$.  With at least two factors one has
$e=5^t\ge25$, contradicting (BI.83).  The only remaining
possibility is a single factor of degree five with simple quotient
$A_6$.  This proves already the assertion for all $e\ge6$,
without using a determinant condition.

We finish the degree-five case.  The two core rows in (BI.80)
are exactly the two degree-five ordinary characters of $A_6$.
One is the deleted natural permutation representation on six
letters: the natural action is doubly transitive, so its deleted
character is irreducible.  The two characters are interchanged
by the diagonal outer automorphism of
$A_6=\operatorname{PSL}_2(9)$; this is also the half-degree pair
in \cite[Lemma~4.4, p.~9]{BI:WhiteSL2}.  Precomposing by that
automorphism leaves the matrix image of $A_6$ unchanged up to
equivalence.  Thus either core image is the deleted natural
$A_6$ image in dimension five.

Its linear normalizer is contained in $\mathbf C^*S_6$ in that
deleted representation.  Indeed the normalizer modulo scalars
injects into $\operatorname{Aut}(A_6)$ and must stabilize the
chosen degree-five character.  Its outer stabilizer has order two:
the field outer automorphism fixes the half-degree characters,
whereas the diagonal and mixed automorphisms exchange them
\cite[Lemmas~4.4 and~4.6, pp.~9--10]{BI:WhiteSL2}.
The natural $S_6$ supplies that stabilizer.  Schur's lemma then
makes every remaining difference scalar.  Since $J$ contains
the core, it lies in the exact subgroup range
\[
                          A_6\subset J\subset\mathbf C^*S_6.
\]
The previously proved D-BI-08 applies to this deleted-standard
inducing representation, with its larger-field allowance and
actual-polarization conclusion.  It excludes a dominant nonsplit
source.  This disposes of the final case and proves the theorem.
\end{proof}

\paragraph{Scope of the computation.}
The input contains all thirty-three tables in (BI.78), not just
the four exceptional entries.  All representations tested have
irreducible core restriction; incompatible outer automorphisms
are not assumed to act.  The finite computation is only the
spectral input to the complete geometric, field and normalizer
proof.  No signature or branch-parameter equality is used to
infer dominance.  Groups with higher Lie rank and ordinary
alternating groups outside the indicated $A_5,A_6$ identifications
remain separate cases.
